\section{Introduction}\label{sec:introduction}

Enterprise users can make sunk service integrations before a provider determines future pricing. Generative-AI deployment is a motivating case: firms may connect data, software, workflows, or task-specific processes before future usage terms are fully settled. The formal model is deliberately more general, however, and contains no primitive unique to AI. Expected tariff architecture can change the installed user composition, while that composition can change the provider's later tariff incentive. This paper studies that reciprocal feedback.

A monopoly provider serves two task classes. Potential tasks decide whether to incur heterogeneous integration costs and then receive a current non-expropriable benefit. The provider later observes installed masses and chooses an anonymous two-part tariff $T(q)=F+pq$. Flat pricing has $p=0$; metered pricing has $p>0$ and requires a real activation cost $\mu$.

The headline results apply on strict $\mathcal R^+$, where the complete rational-expectations candidate interval is first derived from globally optimal continuation play and both classes are then shown to be globally served throughout that interval. Anticipated metering induces a high-use mass $h_M$ satisfying
\[
 h_0<h_M<h_F,
\]
while the provider's gross gain from metering, $\Phi(l,h)$, rises with $h$. Hence
\[
 \mu_M=\Phi(l,h_M)<\Phi(l,h_F)=\mu_F.
\]
Pure metering is self-consistent for $\mu\leq\mu_M$ and pure flat pricing for $\mu\geq\mu_F$. For the strict middle interval $\mu_M<\mu<\mu_F$, neither pure tariff expectation is self-consistent and the certified branch has a unique aggregate mixed resolution generated by actual provider randomization between the two globally optimal tariffs. If integration is architecture-insensitive, the thresholds coincide.

The welfare result is narrower. Holding the installed base fixed, the provider's private metering threshold exceeds the social threshold for the same decentralized allocations by $d h p^*(h)>0$. Once integration responds endogenously, no unconditional welfare ranking is available.

The ingredients themselves are known. Fixed versus usage-linked pricing, commitment and forward-looking adoption, hold-up after sunk investment, mixed pricing choices, and AI/cloud tariff design all have substantial prior literatures \citep{sundararajan2004,hagiu2006,gans2012,mutherswismer2022,wanghu2014,minryu2025,bergemannbonattismolin2025,bergemannwang2025}. We do not claim novelty for the generic sequence ``expectations $\rightarrow$ early participation/investment $\rightarrow$ later pricing.'' The narrower candidate contribution is that the \emph{anticipated tariff architecture} changes a sunk installed composition, that same composition changes the provider's \emph{relative profitability of the two architectures}, and this feedback splits one fixed-state threshold into two expectation-contingent self-consistency thresholds. Section~\ref{sec:literature} states the limits of the current prior-art comparison, including the source limitation for Hagiu's published 2006 version.

The scope is equally narrow. Outside strict $\mathcal R^+$, another continuation active set can dominate. Strict concavity alone does not guarantee the state ordering, so nonquadratic cases are treated only as numerical robustness. The paper therefore makes no global all-parameter, arbitrary-concavity, non-uniform-CDF, or AI-specific theorem claim. Whether the GenAI-forward title or a more general digital-service title is preferable is a Stage-12 journal-positioning question, not a result of the present theory.

The remainder of the paper presents the model, equilibrium, welfare, robustness and scope, institutional interpretation, related literature, discussion, and conclusion.
